% =====================================================================
% Section 7 --- Conclusion (merged draft; supersedes (7)discussion.tex
% and the stub (8)conclusion.tex)
% ---------------------------------------------------------------------
% Job: return to the thesis; position against BF as the frictionless pole;
%      carry the closure lineage (Sen, Taylor, Rose, FIDELIO); state the
%      limitations; close with practical guidance. NO bridge-discovery
%      claims -- the bridge is measurement. Kuhnian frame kept minimal.
% Mindmap: N03, N02D.
% Status: first full draft (2026-10-22). Absorbs the two paragraphs the
%      dropped Discussion would have carried (positioning, limitations).
% NUMBERS TO INSERT once the Sec. 5 prose is final: sign-reversal count and
%      the multiplier range across the five labor closures (Tables 4/5);
%      the Delta Pi/Y range from the financing wedge (Sec. 5.3).
% =====================================================================
\section{Conclusion}\label{sec:conclusion}

This paper began from a suspicion about a genre of numbers. Input--output
models of a socio-ecological transformation report multipliers, employment
effects and fiscal bills as if the short run were a matter of arithmetic:
invert a matrix, read off the change. The suspicion is that the arithmetic
is the last step, and the least interesting one. What precedes it --- the
wage closure and the financing regime --- is a choice among competing
theories of what the short run is, and that choice is not a technicality.
It is first-order for the magnitude of the result, and occasionally for its
sign.

Our contribution is not a new model but a discipline of comparison. One
kernel, five labor-market closures crossed with three financing regimes,
one empirically informed impulse: fifteen scenarios whose differences are
entirely attributable to the assumptions that define them. Read this way,
the matrix does two things at once. It identifies where the closure choice
is redundant --- the pegged real wage and its Leontief fixed-coefficient
corner coincide exactly, so that two rows of the grid collapse into one
economy and the basic intuition of the two-pole picture survives intact.
And it identifies where the choice decides the answer: the row with elastic
sectoral labor supply is a distinct economy, with its own multipliers and
its own distributional profile, and the financing regime shifts the fiscal
incidence even where the labor-market treatment does not.

This result is neither a refutation nor a correction of the
production-network literature; it locates that literature within the family.
The accounting with intermediate inputs of \citet{Hulten.1978} and its
modern descendants \citep{Baqaee.2021,Acemoglu.2012} describe the
frictionless pole of our matrix --- the limit in which sectoral labor supply
is infinitely elastic and reallocation carries no cost. Beyond Hulten, the
question is what the distance from that pole is worth. Our answer is that
the distance is measurable rather than rhetorical, and that it is large
enough to matter for policy evaluation: the bridge between the paradigms is
not a new theory but a common measuring rod.

The five closures are no more an invention than the pole is. They are the
canonical answers of a long debate --- \citet{Sen1963} on the demarcation of
neoclassical and neo-Keynesian closures, \citet{Taylor1990} on structuralist
closures, \citet{Rose.1995} on their translation into computable form ---
carried into an accounting structure that transformation research already
uses. FIDELIO \citep{Rocchi.2025} operationalizes one of these closures; the
matrix makes that closure an explicit modelling decision rather than a
background assumption.

Three limitations bound the claims. First, labor is the only production
factor in the model, so our results speak to the labor-market channel of the
short run and are silent on capital-, capacity- and investment-driven
adjustment; the Kaldor--Verdoorn and markup doors remain open. Second, the
horizon is static: we compare short-run equilibria, not adjustment paths.
Third, the welfare metric is dispersion-blind --- real consumption deflated
by $\mathcal{P}$ says nothing about who within the household sector bears
the adjustment. The skill-class extension requested by the referee is
therefore the natural next step, and the open economy and dynamic closures
are likewise outside the scope of this paper.

For the applied modeller, the matrix is a checklist rather than a verdict.
Name the wage closure. Name the financing regime. Then check whether the
chosen cell is one in which the choice is redundant --- in which case report
it as redundant --- or one in which it decides the answer, in which case
report the range. Reported transformation multipliers are not wrong; they
are conditional, and the condition is usually left implicit. Making it
explicit costs little and changes how the numbers read.

The older lesson survives the newer formalism. The short run is a modelling
choice, and the choice carries the paradigm with it. Which short run --- and
whose money --- remains the question; the numbers answer it only after the
assumptions have been named.
